
\label{sec:robots}
\begin{figure}[h]
    \centering
    \includegraphics[width=\linewidth]{compressed_images_final/pi07robots_v5.pdf}
    \caption{\textbf{Illustrations of some of the robots in our experiments.} We evaluate \MethodName{} on a variety of robots, including bimanual mobile manipulators (left), static bimanual robots (middle), and a bimanual UR5e setup (right) that we use for cross-embodiment experiments.}
    \label{fig:robots}
\vspace{-5pt}
\end{figure}

\begin{figure*}[h!]
    \centering
    \includegraphics[width=\linewidth]{compressed_images_final/full_film_strip.pdf}
    \caption{\textbf{Illustration of selected evaluation tasks.} We evaluate \MethodName{} on a number of tasks, and two of the more longer-horizon ones are visualized here. For some tasks such as ``Take Out Trash'', we provide a coarse instruction like ``take out the trash'' and \MethodName{} performs the full long-horizon task. For other tasks which do not appear in the training data for \MethodName{} such as ``Toasting a Bagel'', we can leverage the strong language following capabilities of \MethodName{} to coach it to perform the task with a series of detailed instructions that break down the task step-by-step.}
    \label{fig:multi_task_film_strip}
\vspace{-6pt}
\end{figure*}

We deploy \MethodName{} in a variety of robot platforms (Fig.~\ref{fig:robots}), including bimanual mobile manipulators with two 6 DoF arms, static bimanual manipulators with lightweight 6 DoF arms (``BiPi''), and a bimanual UR5e system with Robotiq grippers, which we use for cross-embodiment experiments. Additional generalization and language following experiments use a single-arm 6 DoF system, using the same arms as the BiPi platform. Note that while a large fraction of our data is collected with arms that resemble the BiPi platform, the UR5e arms that we use for cross-embodiment testing are significantly longer, have a different morphology, and are much heavier. In practice, the UR5e arms need to employ a different manipulation strategy due to the shape of the arms, their positioning over the table (on the sides rather than at one edge), and the shape of the gripper and fingers, making cross-embodiment transfer to this platform a significant challenge.
All manipulators use parallel-jaw grippers. The UR5e robots run at 20 Hz, while all other robots run at 50 Hz. Each robot has a front-facing camera as well as wrist camera on each arm, and the mobile robots also have a rear-facing camera.
The action output of the \MethodName{} model is applied on each robot using a simple PD controller. For commanding end-effector movement, we apply numerical inverse kinematics to convert target end-effector poses into target joint positions.

\begin{figure*}[h!]
    \centering
    \includegraphics[width=\textwidth]{compressed_images_final/distillation_results.pdf}
    \caption{\textbf{Out-of-the-box dexterity: }\MethodName{} can perform a wide range of highly dexterous tasks directly out of the box. We consider tasks from \Piszs{} \citep{pistar06} (top row) and a number of other dexterous tasks including ones from the ``Robot Olympics'' experiments (bottom row). For the tasks from \Piszs{}, we report success rate and normalized throughput (relative to the specialist model; raw throughput means successes per hour), while for other tasks we report task progress. We find that the same \MethodName{} model can match the performance of the task-specific post-trained specialist policy from \Piszs{} or \Pizs{} for each of these tasks, and even achieve higher throughput than the RL specialists in diverse laundry folding and box building.}
    \label{fig:distillation_results}
\vspace{-10pt}
\end{figure*}

\begin{figure*}[h!]
    \centering
    \includegraphics[width=\textwidth]{compressed_images_final/distillation_ablations.pdf}
    \caption{\textbf{Impact of prompt composition and evaluation data on out-of-the-box performance: }We compare \MethodName{} with two ablations: one that does not include episode metadata in the context, \MethodName{} (no metadata), and another that does not include data from autonomous evaluation episodes during training, \MethodName{} (no eval data). We find that \MethodName{} outperforms both \MethodName{} (no metadata) and  \MethodName{} (no eval data) across the board, with the gap most prominent in throughput. Throughput (successes/hour) here is normalized relative to \MethodName{}.}
    \label{fig:distillation_ablations}
\vspace{-5pt}
\end{figure*}

\begin{figure}[h]
    \centering
    \includegraphics[width=\linewidth]{compressed_images_final/history.pdf}
    \caption{\textbf{Tasks that require memory: }\MethodName{} can also perform tasks that require explicitly keeping track of prior context, achieving similar or better performance compared to the specialist policies with memory fine-tuned to some of the tasks in the MEM paper~\citep{torne2026mem}.}
    \label{fig:memory}
\vspace{-10pt}
\end{figure}

\begin{figure*}[h]
    \centering
    \includegraphics[width=\textwidth]{compressed_images_final/instruction_following.pdf}
    \caption{\textbf{Broad instruction following in novel environments: } We evaluate \MethodName{} on 14 instruction following scenarios, each of which involve following a sequence of 3-6 open-ended instructions, across 4 unseen kitchen and 2 unseen bedroom environments. We report the instruction following success rate, the percentage of total instructions that were correctly followed across all evaluations. We find that \MethodName{} significantly outperforms \Pizf{} and \Pizs{} across the board, attaining high absolute success rates.}
    \label{fig:instruction_following}
\vspace{-8pt}
\end{figure*}

\begin{figure}[h!]
    \centering
    \includegraphics[width=\linewidth]{compressed_images_final/instruction_generalization.pdf}
    \caption{\textbf{Following complex referential instructions: }\MethodName{} and prior models all succeed on the simpler re-arrangement instructions (instructions include ``pick up the spoon'', ``put the spoon to the left of the fork'' and ``put the spoon to the right of the fork''), but \MethodName{} performs significantly better on the complex and unusual instructions (instructions include ``pick up the largest bowl on the table'', ``pick up the object I would use to eat soup'' and ``pick up the fruit on the largest plate''). Including subgoal images generated by a lightweight world model (\MethodName{} (GC)) further boosts instruction following performance, making \MethodName{} significantly more capable at following complex instructions.}\label{fig:instruction_generalization}
\vspace{-3pt}
\end{figure}

\begin{figure}[h!]
    \centering
    \includegraphics[width=\linewidth]{compressed_images_final/compositional_generalization.pdf}
    \caption{\textbf{Breaking dataset biases by following instructions: }the improved language-following performance of \MethodName{} enables it to break strong dataset biases. Prior models struggle on these data bias challenge tasks, which require following instructions that contradict the pattern in the data (e.g., putting dishes in the trash and trash in the dish bin). \MethodName, however, can follow instructions sufficiently well that it can break these strong biases to still perform the task. 
    Notably, for the ``Reverse Fridge to Microwave'' task, including subgoal images from the world model in the context (\MethodName{}(GC)) is critical for success.}
    \label{fig:compositional_generalization}
\vspace{-10pt}
\end{figure}

